\documentclass[10pt,a4paper]{amsart}
\usepackage{amsmath,amssymb,amsthm,graphicx}
% Retained preamble for the Stacks Project chapter divisors.tex (Divisors).
% The chapter file pinned-original-source.tex carries only \input{preamble} on its first line, so no
% preamble text lives inside the chapter itself.  The Stacks Project uses one single shared preamble for
% every chapter; the snapshot stages it verbatim in sources/stacks-chow/preamble.tex.  The block below is
% copied mechanically, line for line, from that file's lines 14-18 (the Xy-pic setup) and lines 154-217
% (the native theorem declarations together with the mathematical macros).  Nothing was added, removed or
% redefined.  The wrapper class and the amsmath/amssymb/amsthm/graphicx base packages are supplied by the
% assembler, so the class/\externaldocument part of the shared preamble is not repeated here.
% For commutative diagrams we use Xy-pic
\usepackage[all]{xy}

% We use 2cell for 2-commutative diagrams.
\xyoption{2cell}
% Theorem environments.
%
\theoremstyle{plain}
\newtheorem{theorem}[subsection]{Theorem}
\newtheorem{proposition}[subsection]{Proposition}
\newtheorem{lemma}[subsection]{Lemma}

\theoremstyle{definition}
\newtheorem{definition}[subsection]{Definition}
\newtheorem{example}[subsection]{Example}
\newtheorem{exercise}[subsection]{Exercise}
\newtheorem{situation}[subsection]{Situation}

\theoremstyle{remark}
\newtheorem{remark}[subsection]{Remark}
\newtheorem{remarks}[subsection]{Remarks}

\numberwithin{equation}{subsection}

% Macros
%
\def\lim{\mathop{\mathrm{lim}}\nolimits}
\def\colim{\mathop{\mathrm{colim}}\nolimits}
\def\Spec{\mathop{\mathrm{Spec}}}
\def\Hom{\mathop{\mathrm{Hom}}\nolimits}
\def\Ext{\mathop{\mathrm{Ext}}\nolimits}
\def\SheafHom{\mathop{\mathcal{H}\!\mathit{om}}\nolimits}
\def\SheafExt{\mathop{\mathcal{E}\!\mathit{xt}}\nolimits}
\def\Sch{\mathit{Sch}}
\def\Mor{\mathop{\mathrm{Mor}}\nolimits}
\def\Ob{\mathop{\mathrm{Ob}}\nolimits}
\def\Sh{\mathop{\mathit{Sh}}\nolimits}
\def\NL{\mathop{N\!L}\nolimits}
\def\CH{\mathop{\mathrm{CH}}\nolimits}
\def\proetale{{pro\text{-}\acute{e}tale}}
\def\etale{{\acute{e}tale}}
\def\QCoh{\mathit{QCoh}}
\def\Ker{\mathop{\mathrm{Ker}}}
\def\Im{\mathop{\mathrm{Im}}}
\def\Coker{\mathop{\mathrm{Coker}}}
\def\Coim{\mathop{\mathrm{Coim}}}

% Boxtimes
%
\DeclareMathSymbol{\boxtimes}{\mathbin}{AMSa}{"02}

%
% Macros for moduli stacks/spaces
%
\def\QCohstack{\mathcal{QC}\!\mathit{oh}}
\def\Cohstack{\mathcal{C}\!\mathit{oh}}
\def\Spacesstack{\mathcal{S}\!\mathit{paces}}
\def\Quotfunctor{\mathrm{Quot}}
\def\Hilbfunctor{\mathrm{Hilb}}
\def\Curvesstack{\mathcal{C}\!\mathit{urves}}
\def\Polarizedstack{\mathcal{P}\!\mathit{olarized}}
\def\Complexesstack{\mathcal{C}\!\mathit{omplexes}}
% \Pic is the operator that assigns to X its picard group, usage \Pic(X)
% \Picardstack_{X/B} denotes the Picard stack of X over B
% \Picardfunctor_{X/B} denotes the Picard functor of X over B
\def\Pic{\mathop{\mathrm{Pic}}\nolimits}
\def\Picardstack{\mathcal{P}\!\mathit{ic}}
\def\Picardfunctor{\mathrm{Pic}}
\def\Deformationcategory{\mathcal{D}\!\mathit{ef}}

% --- appended 2026-10-12 (Chief): the project own definitions of the reference/slogan/history
% display environments, copied mechanically from the shared Stacks preamble lines 8-12
% (\usepackage{verbatim} plus the three \newenvironment lines, which the chapter preamble keeps
% verbatim from the project preamble). Nothing added, removed or redefined: these environments
% carry no mathematical content and the project itself typesets them as comments. Without these
% three lines a counted statement whose own text uses them fails with Environment slogan undefined.
\usepackage{verbatim}
\newenvironment{reference}{\comment}{\endcomment}
\newenvironment{slogan}{\comment}{\endcomment}
\newenvironment{history}{\comment}{\endcomment}
\usepackage{cleveref}
\crefname{originallemma}{original lemma}{original lemmas}
\Crefname{originallemma}{Original lemma}{Original lemmas}
\makeatletter
\newlabel{algebra-lemma-polynomial-ring-UFD}{{0BC1}{source}{}{sourcepointer-1}{}}
\newlabel{algebra-lemma-polynomial-ring-UFD@cref}{{[originallemma][0][]0BC1}{[1][0][]source}{}{}{}}
\makeatother
\crefname{originallemma}{original lemma}{original lemmas}
\Crefname{originallemma}{Original lemma}{Original lemmas}
\makeatletter
\newlabel{coherent-lemma-cohomology-projective-space-over-ring}{{01XT}{source}{}{sourcepointer-2}{}}
\newlabel{coherent-lemma-cohomology-projective-space-over-ring@cref}{{[originallemma][0][]01XT}{[1][0][]source}{}{}{}}
\makeatother
\crefname{originallemma}{original lemma}{original lemmas}
\Crefname{originallemma}{Original lemma}{Original lemmas}
\makeatletter
\newlabel{lemma-Pic-projective-space-UFD}{{0BXJ}{source}{}{sourcepointer-3}{}}
\newlabel{lemma-Pic-projective-space-UFD@cref}{{[originallemma][0][]0BXJ}{[1][0][]source}{}{}{}}
\makeatother
\crefname{originallemma}{original lemma}{original lemmas}
\Crefname{originallemma}{Original lemma}{Original lemmas}
\makeatletter
\newlabel{lemma-codimension-1-is-effective-Cartier}{{0BXH}{source}{}{sourcepointer-4}{}}
\newlabel{lemma-codimension-1-is-effective-Cartier@cref}{{[originallemma][0][]0BXH}{[1][0][]source}{}{}{}}
\makeatother
\title{A closed subscheme of projective space over a Noetherian UFD all of whose components have codimension one is cut out by a single homogeneous polynomial}
\author{The Stacks Project Authors}
\date{stacks-project:divisors.tex:local-snapshot-20261012; GNU FDL 1.2 or later}
\begin{document}
\maketitle
\noindent\textit{Editorial note.} All delivered mathematical text is the author's own, copied from the pinned Stacks Project chapter file pinned-original-source.tex; no statement or proof step was rewritten, reordered or removed.\par\medskip
\noindent\textit{Editorial note.} Following the multi-claim convention for this source, source\_url is the source-level URL https://stacks.math.columbia.edu/ and the source label lemma-equation-codim-1-in-projective-space is carried in claim\_key and here rather than in a per-claim tag URL.\par\medskip
\noindent\textit{Editorial note.} The preamble is the chapter's shared Stacks Project preamble, retained as sources/stacks-divisors/retained-preamble.tex: a line-for-line excerpt of the shared Stacks preamble staged with this snapshot. No macro was invented or redefined, and this card adds no preamble text of its own.\par\medskip
\noindent\textit{Editorial note.} Statement span 8041-8049 and proof span 8051-8091 come from the halfA pool card operations/pool-cards/stacks-divisors-halfA.json and were re-read in the pinned chapter file before this card was written.\par\medskip
\noindent\textit{Editorial note.} The proof displays the dehomogenized equation on each standard open and the divisibility relations after clearing denominators; both displays are preserved unchanged.\par\medskip
\noindent\textit{Editorial note.} The author writes the ideal of the subscheme as the kernel of the map to the sum of the twists of the structure sheaf's sections; that description is kept as written.\par\medskip
\noindent\textit{Editorial note.} 4 cross-references occur inside the counted spans; they are listed in external\_references as navigation pointers only.\par\medskip
\noindent\textit{Allowed standard prerequisites (editorial).}\begin{itemize}
\item The polynomial ring over a unique factorisation domain is a unique factorisation domain, cited from the algebra chapter. Cited at proof lines 8053-8056 and again at 8084-8085. Shows that the local rings of projective space over the base are unique factorisation domains, and later allows the divisibility argument on homogeneous polynomials.
\item A closed subscheme without embedded points all of whose components have codimension one in a scheme with unique factorisation domain local rings is an effective Cartier divisor, cited from the same chapter. Cited at proof lines 8057-8058. Turns the given closed subscheme into an effective Cartier divisor, hence its ideal sheaf is invertible.
\item On projective space over a unique factorisation domain every invertible sheaf is a twist of the structure sheaf, cited from the same chapter. Cited at proof lines 8060-8062. Identifies the invertible ideal sheaf with a twist of the structure sheaf.
\item The description of sections of twists of the structure sheaf on projective space over a ring as homogeneous polynomials of the corresponding degree, cited from the cohomology of schemes chapter. Cited at proof lines 8066-8068. Produces the single homogeneous polynomial cutting out the divisor on each standard open.
\end{itemize}\medskip
\ifdefined\hypertarget\hypertarget{sourcepointer-1}{}\fi
\noindent\textit{Source reference.} Original lemma 0BC1: focused-seeds/AGP/sources/stacks-snapshot-20261011/tags.txt, source lines 14264--14264. Same-chapter reference used in the retained author text at source line(s) 7201, 8055. It is supplied only as a navigation pointer so that the retained text compiles; the counted item delivers the author's own statement and the author's complete proof, and no part of the referenced item is copied in or substituted for the counted argument.\par\medskip
\ifdefined\hypertarget\hypertarget{sourcepointer-2}{}\fi
\noindent\textit{Source reference.} Original lemma 01XT: focused-seeds/AGP/sources/stacks-snapshot-20261011/tags.txt, source lines 2532--2532. Same-chapter reference used in the retained author text at source line(s) 8071. It is supplied only as a navigation pointer so that the retained text compiles; the counted item delivers the author's own statement and the author's complete proof, and no part of the referenced item is copied in or substituted for the counted argument.\par\medskip
\ifdefined\hypertarget\hypertarget{sourcepointer-3}{}\fi
\noindent\textit{Source reference.} Original lemma 0BXJ: focused-seeds/AGP/sources/stacks-snapshot-20261011/tags.txt, source lines 14982--14982. Same-chapter reference used in the retained author text at source line(s) 8062. It is supplied only as a navigation pointer so that the retained text compiles; the counted item delivers the author's own statement and the author's complete proof, and no part of the referenced item is copied in or substituted for the counted argument.\par\medskip
\ifdefined\hypertarget\hypertarget{sourcepointer-4}{}\fi
\noindent\textit{Source reference.} Original lemma 0BXH: focused-seeds/AGP/sources/stacks-snapshot-20261011/tags.txt, source lines 14980--14980. Same-chapter reference used in the retained author text at source line(s) 3187, 8057. It is supplied only as a navigation pointer so that the retained text compiles; the counted item delivers the author's own statement and the author's complete proof, and no part of the referenced item is copied in or substituted for the counted argument.\par\medskip
\begin{lemma}
\label{lemma-equation-codim-1-in-projective-space}
Let $R$ be a Noetherian UFD. Let $Z \subset \mathbf{P}^n_R$ be a
closed subscheme
which has no embedded points such that every irreducible component
of $Z$ has codimension $1$ in $\mathbf{P}^n_R$.
Then the ideal $I(Z) \subset R[T_0, \ldots, T_n]$ corresponding
to $Z$ is principal.
\end{lemma}

\begin{proof}
Observe that the local rings of $X = \mathbf{P}^n_R$ are
UFDs because $X$ is covered by affine pieces isomorphic
to $\mathbf{A}^n_R$ and $R[x_1, \ldots, x_n]$ is a UFD
(Algebra, Lemma \ref{algebra-lemma-polynomial-ring-UFD}).
Thus $Z$ is an effective Cartier divisor by
Lemma \ref{lemma-codimension-1-is-effective-Cartier}.
Let $\mathcal{I} \subset \mathcal{O}_X$ be the quasi-coherent
sheaf of ideals corresponding to $Z$.
Choose an isomorphism $\mathcal{O}(m) \to \mathcal{I}$
for some $m \in \mathbf{Z}$, see
Lemma \ref{lemma-Pic-projective-space-UFD}.
Then the composition
$$
\mathcal{O}_X(m) \to \mathcal{I} \to \mathcal{O}_X
$$
is nonzero. We conclude that $m \leq 0$ and that the corresponding
section of $\mathcal{O}_X(m)^{\otimes -1} = \mathcal{O}_X(-m)$
is given by some $F \in R[T_0, \ldots, T_n]$ of degree $-m$, see
Cohomology of Schemes, Lemma
\ref{coherent-lemma-cohomology-projective-space-over-ring}.
Thus on the $i$th standard open $U_i = D_+(T_i)$ the
closed subscheme $Z \cap U_i$ is cut out by the ideal
$$
(F(T_0/T_i, \ldots, T_n/T_i)) \subset R[T_0/T_i, \ldots, T_n/T_i]
$$
Thus the homogeneous elements of the graded ideal
$I(Z) = \Ker(R[T_0, \ldots, T_n] \to \bigoplus \Gamma(\mathcal{O}_Z(m)))$
is the set of homogeneous polynomials $G$ such that
$$
G(T_0/T_i, \ldots, T_n/T_i) \in (F(T_0/T_i, \ldots, T_n/T_i))
$$
for $i = 0, \ldots, n$. Clearing denominators, we see there exist
$e_i \geq 0$ such that
$$
T_i^{e_i}G \in (F)
$$
for $i = 0, \ldots, n$. As $R$ is a UFD, so is $R[T_0, \ldots, T_n]$.
Then $F | T_0^{e_0}G$ and $F | T_1^{e_1}G$ implies $F | G$ as
$T_0^{e_0}$ and $T_1^{e_1}$ have no factor in common. Thus $I(Z) = (F)$.
\end{proof}

\end{document}
